%%%%%%%%%%%%%%%%%%%%%%%%%%%%%%%%%%%%%%%%%%%%%%%%%%%%%%%%%%%%
%  Secure Communication Suite – Project Report
%  CSE451: Computer and Network Security
%  Ain Shams University – Faculty of Engineering
%  CHEP – Spring 2026
%%%%%%%%%%%%%%%%%%%%%%%%%%%%%%%%%%%%%%%%%%%%%%%%%%%%%%%%%%%%
\documentclass[12pt,a4paper]{article}

% ── Packages ─────────────────────────────────────────────
\usepackage[utf8]{inputenc}
\usepackage[T1]{fontenc}
\usepackage{lmodern}
\usepackage[margin=2.5cm]{geometry}
\usepackage{graphicx}
\usepackage{xcolor}
\usepackage{hyperref}
\usepackage{listings}
\usepackage{enumitem}
\usepackage{titlesec}
\usepackage{fancyhdr}
\usepackage{amsmath}
\usepackage{booktabs}
\usepackage{caption}
\usepackage{float}

% ── Colours ──────────────────────────────────────────────
\definecolor{codegreen}{rgb}{0.25,0.49,0.37}
\definecolor{codegray}{rgb}{0.5,0.5,0.5}
\definecolor{codepurple}{rgb}{0.58,0,0.82}
\definecolor{backcolour}{rgb}{0.97,0.97,0.97}
\definecolor{asublue}{RGB}{0,70,140}

% ── Listings Style ───────────────────────────────────────
\lstdefinestyle{pythonstyle}{
    backgroundcolor=\color{backcolour},
    commentstyle=\color{codegreen},
    keywordstyle=\color{blue}\bfseries,
    numberstyle=\tiny\color{codegray},
    stringstyle=\color{codepurple},
    basicstyle=\ttfamily\footnotesize,
    breakatwhitespace=false,
    breaklines=true,
    captionpos=b,
    keepspaces=true,
    numbers=left,
    numbersep=8pt,
    showspaces=false,
    showstringspaces=false,
    showtabs=false,
    tabsize=4,
    language=Python,
    frame=single,
    rulecolor=\color{black!30},
    xleftmargin=1.5em,
    framexleftmargin=1em,
}
\lstset{style=pythonstyle}

% ── Header / Footer ─────────────────────────────────────
\pagestyle{fancy}
\fancyhf{}
\rhead{\small CSE451 – Spring 2026}
\lhead{\small Secure Communication Suite}
\cfoot{\thepage}

% ── Hyperref setup ───────────────────────────────────────
\hypersetup{
    colorlinks=true,
    linkcolor=asublue,
    urlcolor=asublue,
    citecolor=asublue,
}

% ── Title-page helpers ───────────────────────────────────
\newcommand{\HRule}{\rule{\linewidth}{0.5mm}}

%%%%%%%%%%%%%%%%%%%%%%%%%%%%%%%%%%%%%%%%%%%%%%%%%%%%%%%%%%%%
\begin{document}

% ===========================================================
%  TITLE PAGE
% ===========================================================
\begin{titlepage}
\centering
\vspace*{1cm}
{\large\textbf{Ain Shams University}}\\[0.3cm]
{\large Faculty of Engineering}\\[0.1cm]
{\large Computer and Systems Engineering Department}\\[0.1cm]
{\large CHEP -- Spring 2026}\\[1.5cm]
\HRule\\[0.6cm]
{\Huge\bfseries Secure Communication Suite}\\[0.3cm]
{\Large A Cryptography Application in Python}\\[0.4cm]
\HRule\\[1.5cm]
{\large\textbf{CSE451: Computer and Network Security}}\\[2cm]

% ── Group Members (fill in your names) ───────────────────
\begin{minipage}{0.5\textwidth}
\begin{flushleft}\large
\textbf{Group Members:}\\[0.3cm]
Marwan Noureldin – 21P0165\\
Yehia Mahmoud – 21P0166\\
Ahmed Ayman – 21P002\\
Youssef Ayman – 21P0010 \\
Hanna Mohamed – 21P0162\\
\end{flushleft}
\end{minipage}
\vfill
{\large April 2026}
\end{titlepage}

% ===========================================================
%  TABLE OF CONTENTS
% ===========================================================
\tableofcontents
\newpage

% ===========================================================
%  PHASE 2 – Development of Cryptographic Modules
% ===========================================================
\section{Phase 2: Development of Cryptographic Modules}

This phase focuses on implementing the three fundamental cryptographic building blocks of the Secure Communication Suite: the \textbf{Block Cipher Module} (AES), the \textbf{Public Key Cryptosystem Module} (RSA), and the \textbf{Hashing Module} (SHA-256). Each module is encapsulated in its own Python class inside the \texttt{src/crypto/} package.

% -----------------------------------------------------------
\subsection{Block Cipher Module – AES (\texttt{block\_cipher.py})}
% -----------------------------------------------------------

\subsubsection{Overview}
The Advanced Encryption Standard (AES) is a symmetric-key block cipher adopted as the encryption standard by the U.S.\ National Institute of Standards and Technology (NIST). It operates on fixed-size blocks of 128 bits and supports key sizes of 128, 192, or 256 bits.

Our implementation uses \textbf{AES-128} (16-byte key) with the \textbf{EAX mode of operation}. EAX is an authenticated encryption mode, meaning that it simultaneously provides both \emph{confidentiality} and \emph{integrity/authenticity}. This is achieved by producing a \emph{nonce} and an authentication \emph{tag} alongside the ciphertext.

\subsubsection{Why EAX Mode?}
\begin{itemize}
    \item \textbf{Authenticated Encryption:} Unlike ECB or CBC modes, EAX guarantees that any tampering with the ciphertext will be detected during decryption, because the tag verification will fail.
    \item \textbf{No IV Management Burden:} EAX automatically generates and manages a unique nonce for each encryption call, preventing nonce-reuse vulnerabilities.
    \item \textbf{Stream-Friendly:} EAX can encrypt data of arbitrary length; no padding is needed, eliminating padding-oracle attack vectors.
\end{itemize}

\subsubsection{Implementation Details}

\begin{lstlisting}[caption={AES Block Cipher Module – \texttt{block\_cipher.py}}]
from Crypto.Cipher import AES
from Crypto.Random import get_random_bytes

class AESCipher:
    """
    AES encryption module using EAX mode.
    """
    def __init__(self, key: bytes = None):
        if key is None:
            self.key = get_random_bytes(16)  # AES-128
        else:
            if len(key) not in [16, 24, 32]:
                raise ValueError("Key must be 16, 24, or 32 bytes")
            self.key = key

    def encrypt(self, plaintext: bytes) -> bytes:
        """
        Encrypts plaintext using AES-EAX mode.
        Returns nonce + tag + ciphertext.
        """
        cipher = AES.new(self.key, AES.MODE_EAX)
        nonce = cipher.nonce
        ciphertext, tag = cipher.encrypt_and_digest(plaintext)
        return nonce + tag + ciphertext

    def decrypt(self, encrypted_data: bytes) -> bytes:
        """
        Decrypts data encrypted by the encrypt method.
        """
        nonce = encrypted_data[:16]
        tag = encrypted_data[16:32]
        ciphertext = encrypted_data[32:]
        cipher = AES.new(self.key, AES.MODE_EAX, nonce=nonce)
        plaintext = cipher.decrypt_and_verify(ciphertext, tag)
        return plaintext
\end{lstlisting}

\subsubsection{How the Block Cipher Module Works}

\begin{enumerate}
    \item \textbf{Key Initialization:} The constructor accepts an optional key parameter. If no key is provided, a cryptographically secure random 16-byte key is generated using \texttt{get\_random\_bytes()}. If a key is supplied, its length is validated to be 16, 24, or 32 bytes (corresponding to AES-128, AES-192, or AES-256 respectively).
    
    \item \textbf{Encryption Process:}
    \begin{enumerate}
        \item A new AES cipher object is created in EAX mode, which internally generates a fresh 16-byte \textbf{nonce}.
        \item The plaintext is encrypted and a 16-byte \textbf{authentication tag} is computed via \texttt{encrypt\_and\_digest()}.
        \item The output is a single byte string structured as: \texttt{nonce (16\,B) $\|$ tag (16\,B) $\|$ ciphertext}.
    \end{enumerate}
    
    \item \textbf{Decryption Process:}
    \begin{enumerate}
        \item The received byte string is split into nonce (first 16 bytes), tag (next 16 bytes), and ciphertext (remainder).
        \item A new AES cipher is initialized with the \emph{same key} and the extracted nonce.
        \item \texttt{decrypt\_and\_verify()} recovers the plaintext and simultaneously verifies the tag. If the data has been tampered with, a \texttt{ValueError} is raised.
    \end{enumerate}
\end{enumerate}

\subsubsection{Data Format}
The encrypted output follows a fixed layout, as summarized in Table~\ref{tab:aes-format}.

\begin{table}[H]
\centering
\caption{AES-EAX encrypted data format}
\label{tab:aes-format}
\begin{tabular}{@{}lcc@{}}
\toprule
\textbf{Field} & \textbf{Size (bytes)} & \textbf{Purpose} \\
\midrule
Nonce   & 16 & Ensure unique ciphertexts \\
Tag     & 16 & Authenticate ciphertext \\
Ciphertext & Variable & Encrypted payload \\
\bottomrule
\end{tabular}
\end{table}

% -----------------------------------------------------------
\subsection{Public Key Cryptosystem Module – RSA (\texttt{public\_key.py})}
% -----------------------------------------------------------

\subsubsection{Overview}
RSA (Rivest--Shamir--Adleman) is the most widely deployed public-key cryptosystem. Its security relies on the computational difficulty of factoring the product of two large prime numbers. In our suite, RSA is used specifically for \textbf{secure key exchange}: the client encrypts a randomly generated AES session key with the server's RSA public key, and only the server (which possesses the corresponding private key) can decrypt it.

\subsubsection{PKCS\#1 OAEP Padding}
Our implementation uses \textbf{PKCS\#1 OAEP} (Optimal Asymmetric Encryption Padding), which is the recommended RSA encryption scheme. OAEP provides:
\begin{itemize}
    \item \textbf{Semantic Security:} Identical plaintext encrypted multiple times produces different ciphertext, thwarting chosen-plaintext attacks.
    \item \textbf{Integrity:} The padding scheme includes a hash-based verification that detects malformed ciphertexts.
\end{itemize}

\subsubsection{Implementation Details}

\begin{lstlisting}[caption={RSA Public Key Module – \texttt{public\_key.py}}]
from Crypto.PublicKey import RSA
from Crypto.Cipher import PKCS1_OAEP

class RSACipher:
    """
    RSA asymmetric encryption module for secure key sharing.
    """
    @staticmethod
    def generate_keys(key_size=2048):
        """Generates an RSA keypair (private, public) as bytes."""
        key = RSA.generate(key_size)
        private_key = key.export_key()
        public_key = key.publickey().export_key()
        return private_key, public_key

    @staticmethod
    def encrypt(public_key: bytes, message: bytes) -> bytes:
        """Encrypts a message using a public RSA key."""
        rsa_key = RSA.import_key(public_key)
        cipher = PKCS1_OAEP.new(rsa_key)
        return cipher.encrypt(message)

    @staticmethod
    def decrypt(private_key: bytes, ciphertext: bytes) -> bytes:
        """Decrypts ciphertext using a private RSA key."""
        rsa_key = RSA.import_key(private_key)
        cipher = PKCS1_OAEP.new(rsa_key)
        return cipher.decrypt(ciphertext)
\end{lstlisting}

\subsubsection{Role of the Public Key Cryptosystem Module}

The RSA module fulfils two critical roles in the Secure Communication Suite:

\begin{enumerate}
    \item \textbf{Key Generation:} The \texttt{generate\_keys()} method produces a 2048-bit RSA key pair. The private key is stored securely on the server, while the public key is transmitted to every connecting client. The 2048-bit key size provides a security level of approximately 112 bits, which is considered adequate for current use.
    
    \item \textbf{Hybrid Encryption Bridge:} Pure RSA encryption is limited to small messages (at most 245 bytes with 2048-bit keys and OAEP). Therefore, RSA is used only to securely transmit a 16-byte AES session key from the client to the server, after which all bulk data encryption is handled by the much faster AES algorithm. This is the classical \emph{hybrid encryption} pattern.
    
    \item \textbf{Static Methods:} All methods are \texttt{@staticmethod}, which means the class is used as a stateless utility. Keys are passed explicitly, improving testability and allowing integration with the \texttt{KeyManager}.
\end{enumerate}

% -----------------------------------------------------------
\subsection{Hashing Module – SHA-256 (\texttt{hashing.py})}
% -----------------------------------------------------------

\subsubsection{Overview}
SHA-256 (Secure Hash Algorithm, 256-bit) is a member of the SHA-2 family designed by the NSA. It produces a fixed-size 256-bit (32-byte) digest from an arbitrary-length input. The algorithm guarantees:
\begin{itemize}
    \item \textbf{Pre-image Resistance:} Given a hash $h$, it is computationally infeasible to find any message $m$ such that $\text{SHA-256}(m) = h$.
    \item \textbf{Second Pre-image Resistance:} Given a message $m_1$, it is infeasible to find a different $m_2$ where $\text{SHA-256}(m_1) = \text{SHA-256}(m_2)$.
    \item \textbf{Collision Resistance:} It is infeasible to find any two messages $m_1 \neq m_2$ with the same hash.
\end{itemize}

\subsubsection{Implementation Details}

\begin{lstlisting}[caption={SHA-256 Hashing Module – \texttt{hashing.py}}]
from Crypto.Hash import SHA256

class DeepHash:
    """
    Data integrity checking using SHA-256.
    """
    @staticmethod
    def hash_data(data: bytes) -> str:
        """Computes the SHA-256 hash of the given data."""
        h = SHA256.new()
        h.update(data)
        return h.hexdigest()

    @staticmethod
    def verify(data: bytes, expected_hash: str) -> bool:
        """Verifies SHA-256 hash of data matches expected hash."""
        return DeepHash.hash_data(data) == expected_hash
\end{lstlisting}

\subsubsection{How the Hashing Module Ensures Data Integrity}

\begin{enumerate}
    \item \textbf{Digest Computation – \texttt{hash\_data()}:} The method instantiates a SHA-256 hasher, feeds in the binary data via \texttt{update()}, and returns the hexadecimal-encoded digest string (64 hex characters).
    
    \item \textbf{Verification – \texttt{verify()}:} To check whether data has been altered, the caller provides the data and the previously known hash. The method recomputes the hash and performs a simple equality comparison. If the hashes differ, the data has been tampered with.
    
    \item \textbf{Usage in the Suite:}
    \begin{itemize}
        \item \textbf{Password Storage:} The \texttt{AuthManager} (Phase~3) uses SHA-256 with a random salt to hash passwords before storing them. This ensures that even if the user database is compromised, plaintext passwords are not revealed.
        \item \textbf{GUI Integrity Tool:} The graphical interface provides a ``Data Integrity'' view where users can manually compute SHA-256 digests of arbitrary text, enabling them to verify file or message integrity out-of-band.
        \item \textbf{AES-EAX Tag:} The AES-EAX authentication tag (Phase~2) internally relies on CMAC, but the SHA-256 module serves as an independent integrity layer available to the user.
    \end{itemize}
\end{enumerate}

\subsubsection{Phase 2 Summary}
By the end of Phase~2, the three core cryptographic modules are complete:
\begin{itemize}
    \item \textbf{AES-EAX} provides fast, authenticated symmetric encryption.
    \item \textbf{RSA-OAEP} provides secure asymmetric key exchange.
    \item \textbf{SHA-256} provides tamper-evident hashing.
\end{itemize}

\newpage
% ===========================================================
%  PHASE 3 – Key Management and Authentication Modules
% ===========================================================
\section{Phase 3: Development of Key Management and Authentication Modules}

Phase~3 builds the trust infrastructure that underpins the entire suite. The \textbf{Key Management Module} handles generation, storage, and distribution of RSA key pairs, while the \textbf{Authentication Module} implements password-based user authentication with salted hashing. Both modules reside in the \texttt{src/auth/} package.

% -----------------------------------------------------------
\subsection{Key Management Module (\texttt{key\_management.py})}
% -----------------------------------------------------------

\subsubsection{Overview}
Cryptographic keys are the most sensitive asset in any security system. If a private key is compromised, all data encrypted with the corresponding public key can be decrypted. The \texttt{KeyManager} class addresses three concerns:
\begin{enumerate}
    \item \textbf{Generation:} Creating RSA key pairs when they do not yet exist.
    \item \textbf{Storage:} Persisting keys to the filesystem in PEM format.
    \item \textbf{Distribution:} Sharing public keys between communicating parties.
\end{enumerate}

\subsubsection{Implementation Details}

\begin{lstlisting}[caption={Key Management Module – \texttt{key\_management.py}}]
import os
from crypto.public_key import RSACipher

class KeyManager:
    """Secure methods for key generation, distribution, and storage."""

    @staticmethod
    def get_or_create_keys(user_id: str, key_dir="keys"):
        if not os.path.exists(key_dir):
            os.makedirs(key_dir)

        priv_path = os.path.join(key_dir, f"{user_id}_private.pem")
        pub_path  = os.path.join(key_dir, f"{user_id}_public.pem")

        if os.path.exists(priv_path) and os.path.exists(pub_path):
            with open(priv_path, "rb") as f:
                priv = f.read()
            with open(pub_path, "rb") as f:
                pub = f.read()
            return priv, pub

        priv, pub = RSACipher.generate_keys()

        with open(priv_path, "wb") as f:
            f.write(priv)
        with open(pub_path, "wb") as f:
            f.write(pub)

        return priv, pub

    @staticmethod
    def save_partner_key(partner_id: str, pub_key: bytes,
                         key_dir="keys"):
        if not os.path.exists(key_dir):
            os.makedirs(key_dir)
        path = os.path.join(key_dir, f"{partner_id}_public.pem")
        with open(path, "wb") as f:
            f.write(pub_key)

    @staticmethod
    def load_partner_key(partner_id: str,
                         key_dir="keys") -> bytes:
        path = os.path.join(key_dir, f"{partner_id}_public.pem")
        if not os.path.exists(path):
            return None
        with open(path, "rb") as f:
            return f.read()
\end{lstlisting}

\subsubsection{How the Key Management Module Secures Key Distribution and Storage}

\begin{enumerate}
    \item \textbf{Lazy Key Generation (\texttt{get\_or\_create\_keys}):}
    \begin{itemize}
        \item On first invocation for a given \texttt{user\_id}, the method generates a fresh 2048-bit RSA key pair via \texttt{RSACipher.generate\_keys()} and persists both the private and public key to disk in PEM format.
        \item On subsequent invocations, the method reads the existing keys from disk, avoiding redundant key generation and ensuring key continuity across server restarts.
    \end{itemize}
    
    \item \textbf{Key Storage:}
    \begin{itemize}
        \item All keys are stored in a dedicated \texttt{keys/} directory, created automatically if absent.
        \item Naming convention: \texttt{<user\_id>\_private.pem} and \texttt{<user\_id>\_public.pem}, making key identification straightforward.
        \item Keys are stored in the standard \textbf{PEM} (Privacy Enhanced Mail) format, which is Base64-encoded DER wrapped with header and footer lines. This format is human-readable and widely interoperable.
    \end{itemize}
    
    \item \textbf{Partner Key Distribution (\texttt{save\_partner\_key} / \texttt{load\_partner\_key}):}
    \begin{itemize}
        \item When a client connects, the server sends its public key. The client calls \texttt{save\_partner\_key("server", pub\_key)} to persist this public key locally.
        \item \texttt{load\_partner\_key()} retrieves a previously saved partner's public key. This enables subsequent sessions to verify the server's identity.
    \end{itemize}
    
    \item \textbf{Security Considerations:}
    \begin{itemize}
        \item Private keys remain exclusively on the machine that generated them; they are never transmitted over the network.
        \item In a production environment, the \texttt{keys/} directory should be protected with strict file-system permissions (e.g., \texttt{chmod 700}).
    \end{itemize}
\end{enumerate}

% -----------------------------------------------------------
\subsection{Authentication Module (\texttt{authentication.py})}
% -----------------------------------------------------------

\subsubsection{Overview}
The authentication module implements a \textbf{password-based authentication} system that securely stores credentials and verifies user identities. It uses \textbf{salted SHA-256 hashing} to protect stored passwords.

\subsubsection{Implementation Details}

\begin{lstlisting}[caption={Authentication Module – \texttt{authentication.py}}]
from Crypto.Hash import SHA256
from Crypto.Random import get_random_bytes
import json
import os

class AuthManager:
    """Implement password-based authentication mechanisms."""
    def __init__(self, db_path="users.json"):
        self.db_path = db_path
        self._load_users()

    def _load_users(self):
        if not os.path.exists(self.db_path):
            self.users = {}
            self._save_users()
        else:
            with open(self.db_path, "r") as f:
                self.users = json.load(f)

    def _save_users(self):
        with open(self.db_path, "w") as f:
            json.dump(self.users, f)

    def _hash_password(self, password: str, salt: str) -> str:
        h = SHA256.new()
        h.update((password + salt).encode('utf-8'))
        return h.hexdigest()

    def register(self, username: str, password: str) -> bool:
        if username in self.users:
            return False
        salt = get_random_bytes(16).hex()
        hashed = self._hash_password(password, salt)
        self.users[username] = {"salt": salt, "hash": hashed}
        self._save_users()
        return True

    def authenticate(self, username: str, password: str) -> bool:
        if username not in self.users:
            return False
        user_data = self.users[username]
        hashed = self._hash_password(password, user_data["salt"])
        return hashed == user_data["hash"]
\end{lstlisting}

\subsubsection{Authentication Mechanisms}

The module provides two operations:

\paragraph{Registration.}
\begin{enumerate}
    \item The user supplies a \texttt{username} and \texttt{password}.
    \item The system checks whether the username already exists; if so, registration is rejected to prevent duplicates.
    \item A cryptographically secure random \textbf{salt} (16 bytes, hex-encoded to 32 characters) is generated using \texttt{get\_random\_bytes(16)}.
    \item The password is concatenated with the salt, and the result is hashed with SHA-256: $h = \text{SHA-256}(\texttt{password} \| \texttt{salt})$.
    \item The user record \texttt{\{"salt": ..., "hash": ...\}} is stored in a JSON file (\texttt{users.json}).
\end{enumerate}

\paragraph{Login.}
\begin{enumerate}
    \item The user provides \texttt{username} and \texttt{password}.
    \item The system retrieves the stored salt for that username.
    \item It recomputes $h' = \text{SHA-256}(\texttt{password} \| \texttt{salt})$ and compares $h'$ with the stored hash.
    \item If they match, the user is authenticated; otherwise, login fails.
\end{enumerate}

\subsubsection{How the Authentication Module Verifies User Identities}

\begin{itemize}
    \item \textbf{Salted Hashing Prevents Rainbow-Table Attacks:} Because each user has a unique random salt, precomputed hash tables (rainbow tables) cannot be used to reverse the hash.
    
    \item \textbf{No Plaintext Storage:} Passwords are \emph{never} stored in plaintext. Even if an attacker gains access to \texttt{users.json}, they obtain only salts and hashes, not usable passwords.
    
    \item \textbf{Duplicate Session Prevention:} The server (\texttt{server.py}) maintains a dictionary \texttt{self.clients} mapping connected sockets to usernames. Before authenticating a login, the server checks whether the username already appears in the dictionary. If so, it sends a \texttt{"DUP"} response and refuses the connection, preventing simultaneous duplicate sessions.
    
    \item \textbf{Persistence:} The JSON-based user database is loaded at server startup and updated transactionally on each registration, ensuring data survives server restarts.
\end{itemize}

\subsubsection{Phase 3 Summary}
By the end of Phase~3:
\begin{itemize}
    \item The \textbf{KeyManager} handles RSA key lifecycle (generate, store, load, distribute).
    \item The \textbf{AuthManager} provides secure registration and login with salted SHA-256 password hashing.
    \item Together, these modules provide the trust foundation for Phase~4's integrated system.
\end{itemize}

\newpage
% ===========================================================
%  PHASE 4 – Integration and Testing
% ===========================================================
\section{Phase 4: Integration and Testing}

Phase~4 brings all modules together into a fully functional \textbf{Secure Communication Suite}. This phase covers the \texttt{server.py} (multi-client secure server), \texttt{client.py} (client-side networking), and \texttt{gui.py} (graphical user interface), as well as the testing methodology.

% -----------------------------------------------------------
\subsection{System Architecture}
% -----------------------------------------------------------

The suite follows a classic \textbf{client--server} architecture augmented with a hybrid cryptographic protocol:

\begin{enumerate}
    \item \textbf{Server} (\texttt{server.py}): Listens for incoming TCP connections, authenticates users, performs RSA key exchange, and relays AES-encrypted messages between clients.
    \item \textbf{Client} (\texttt{client.py}): Connects to the server, performs authentication and key exchange, then sends/receives AES-encrypted messages.
    \item \textbf{GUI} (\texttt{gui.py}): Provides a modern graphical front-end built with CustomTkinter, featuring login/register screens, secure chat, and cryptographic tools.
\end{enumerate}

\begin{figure}[H]
\centering
\fbox{\parbox{0.85\textwidth}{\centering
\textbf{Communication Protocol Flow}\\[0.5em]
\texttt{Client} $\xrightarrow{\text{1. LOGIN/REGISTER + credentials}}$ \texttt{Server}\\[0.3em]
\texttt{Client} $\xleftarrow{\text{2. OK / ERR / DUP}}$ \texttt{Server}\\[0.3em]
\texttt{Client} $\xleftarrow{\text{3. RSA Public Key}}$ \texttt{Server}\\[0.3em]
\texttt{Client} $\xrightarrow{\text{4. RSA-encrypted AES key}}$ \texttt{Server}\\[0.3em]
\texttt{Client} $\xleftarrow{\text{5. READY}}$ \texttt{Server}\\[0.3em]
\texttt{Client} $\xleftrightarrow{\text{6. AES-EAX encrypted messages}}$ \texttt{Server}
}}
\caption{Protocol flow for secure session establishment}
\end{figure}

% -----------------------------------------------------------
\subsection{Server Module (\texttt{server.py})}
% -----------------------------------------------------------

\subsubsection{Implementation Details}

\begin{lstlisting}[caption={Secure Server – \texttt{server.py}}]
import socket
import threading
from auth.authentication import AuthManager
from auth.key_management import KeyManager
from crypto.public_key import RSACipher
from crypto.block_cipher import AESCipher

class SecureServer:
    def __init__(self, host='127.0.0.1', port=65432):
        self.host = host
        self.port = port
        self.auth = AuthManager()
        self.server_socket = socket.socket(
            socket.AF_INET, socket.SOCK_STREAM)
        self.server_socket.setsockopt(
            socket.SOL_SOCKET, socket.SO_REUSEADDR, 1)
        self.server_socket.bind((self.host, self.port))
        self.priv_key, self.pub_key = \
            KeyManager.get_or_create_keys("server")
        self.clients = {}            # {socket: username}
        self.broadcast_ciphers = {}  # {socket: AESCipher}

    def start(self):
        self.server_socket.listen()
        print(f"[*] Server listening on {self.host}:{self.port}")
        try:
            while True:
                conn, addr = self.server_socket.accept()
                thread = threading.Thread(
                    target=self.handle_client, args=(conn, addr))
                thread.start()
        except KeyboardInterrupt:
            print("[*] Server shutting down")
            self.server_socket.close()

    def broadcast(self, message_bytes: bytes, sender_conn):
        for client_conn, cipher in \
                self.broadcast_ciphers.items():
            if client_conn != sender_conn:
                try:
                    encrypted = cipher.encrypt(message_bytes)
                    client_conn.send(encrypted)
                except Exception as e:
                    print(f"[-] Broadcast error: {e}")

    def handle_client(self, conn, addr):
        print(f"[*] Connection accepted from {addr}")
        try:
            # 1. Authentication
            auth_msg = conn.recv(1024).decode()
            action, username, password = \
                auth_msg.split('||')

            if action == 'REGISTER':
                success = self.auth.register(
                    username, password)
                conn.send(b"OK" if success else b"ERR")
                if not success: return
            elif action == 'LOGIN':
                if username in self.clients.values():
                    conn.send(b"DUP")
                    return
                success = self.auth.authenticate(
                    username, password)
                conn.send(b"OK" if success else b"ERR")
                if not success: return
            else:
                conn.send(b"ERR")
                return

            # 2. Key Exchange
            conn.send(self.pub_key)
            encrypted_aes_key = conn.recv(1024)
            aes_key = RSACipher.decrypt(
                self.priv_key, encrypted_aes_key)
            cipher = AESCipher(aes_key)

            self.clients[conn] = username
            self.broadcast_ciphers[conn] = cipher
            conn.send(b"READY")

            self.broadcast(
                f"[SERVER] {username} joined!".encode(),
                conn)

            # 3. Secure Communication Loop
            while True:
                encrypted_data = conn.recv(4096)
                if not encrypted_data:
                    break
                plaintext = cipher.decrypt(encrypted_data)
                formatted = \
                    f"[{username}]: {plaintext.decode()}"
                self.broadcast(formatted.encode(), conn)

        except Exception as e:
            print(f"[-] Error handling {addr}: {e}")
        finally:
            username = self.clients.get(conn, "Unknown")
            self.clients.pop(conn, None)
            self.broadcast_ciphers.pop(conn, None)
            conn.close()
            self.broadcast(
                f"[SERVER] {username} left.".encode(), None)

if __name__ == '__main__':
    server = SecureServer()
    server.start()
\end{lstlisting}

\subsubsection{How Modules Are Integrated in the Server}

\begin{enumerate}
    \item \textbf{AuthManager Integration:} At startup, an \texttt{AuthManager} instance is created, loading the user database. Each incoming connection is first processed through \texttt{register()} or \texttt{authenticate()}.
    
    \item \textbf{KeyManager Integration:} On initialization, the server calls \texttt{KeyManager.get\_or\_create\_keys("server")} to obtain its RSA key pair. The public key is sent to each client during the handshake.
    
    \item \textbf{RSACipher Integration:} After the client receives the server's public key, it encrypts a freshly generated AES session key and sends it back. The server uses \texttt{RSACipher.decrypt()} to recover the AES key.
    
    \item \textbf{AESCipher Integration:} A per-client \texttt{AESCipher} object is created with the negotiated session key. All subsequent messages are encrypted and decrypted through this cipher.
    
    \item \textbf{Multi-threaded Design:} Each client connection is handled in a separate thread, enabling concurrent secure sessions. The \texttt{broadcast()} method iterates over all connected clients (except the sender) and encrypts the message with each client's individual AES cipher.
\end{enumerate}

% -----------------------------------------------------------
\subsection{Client Module (\texttt{client.py})}
% -----------------------------------------------------------

\begin{lstlisting}[caption={Secure Client – \texttt{client.py}}]
import socket
import threading
from auth.key_management import KeyManager
from crypto.public_key import RSACipher
from crypto.block_cipher import AESCipher
from Crypto.Random import get_random_bytes

class SecureClient:
    def __init__(self, host='127.0.0.1', port=65432):
        self.host = host
        self.port = port
        self.socket = None
        self.aes_cipher = None
        self.connected = False
        self.username = None

    def connect_and_auth(self, action, username, password):
        try:
            self.socket = socket.socket(
                socket.AF_INET, socket.SOCK_STREAM)
            self.socket.connect((self.host, self.port))

            # 1. Auth Phase
            msg = f"{action}||{username}||{password}"
            self.socket.send(msg.encode())
            res = self.socket.recv(1024).decode()
            if res == "DUP":
                self.socket.close()
                return False, "User already logged in."
            if res != "OK":
                self.socket.close()
                return False, "Authentication failed."

            # 2. Key Exchange Phase
            server_pub_key = self.socket.recv(2048)
            KeyManager.save_partner_key(
                "server", server_pub_key)

            aes_key = get_random_bytes(16)
            encrypted_aes_key = RSACipher.encrypt(
                server_pub_key, aes_key)
            self.socket.send(encrypted_aes_key)

            ready = self.socket.recv(1024)
            if ready == b"READY":
                self.aes_cipher = AESCipher(aes_key)
                self.connected = True
                self.username = username
                return True, "Secure session established."

            return False, "Session setup failed."
        except Exception as e:
            if self.socket:
                self.socket.close()
            return False, f"Connection error: {e}"

    def send_message(self, message: str) -> bool:
        if not self.connected:
            return False
        try:
            encrypted = self.aes_cipher.encrypt(
                message.encode())
            self.socket.send(encrypted)
            return True
        except:
            self.connected = False
            return False

    def receive_message(self) -> str:
        if not self.connected:
            return None
        try:
            encrypted = self.socket.recv(4096)
            if not encrypted:
                self.connected = False
                return None
            plaintext = self.aes_cipher.decrypt(encrypted)
            return plaintext.decode()
        except:
            self.connected = False
            return None

    def close(self):
        if self.socket:
            self.socket.close()
        self.connected = False
\end{lstlisting}

% -----------------------------------------------------------
\subsection{Graphical User Interface (\texttt{gui.py})}
% -----------------------------------------------------------

The GUI is built using \textbf{CustomTkinter}, a modern Python UI library that extends tkinter with contemporary styling. The interface consists of two main screens:

\begin{enumerate}
    \item \textbf{Authentication Screen:} A centred login card with username/password fields and Login/Register buttons.
    \item \textbf{Main Interface:} A sidebar-based layout with:
    \begin{itemize}
        \item \textbf{Secure Chat} – real-time encrypted messaging with a text display, message entry, and send button.
        \item \textbf{Data Integrity (Crypto Tools)} – an SHA-256 hasher and a standalone AES encrypt/decrypt tool.
        \item \textbf{Disconnect} – graceful session teardown.
    \end{itemize}
\end{enumerate}

\begin{lstlisting}[caption={GUI Application – \texttt{gui.py} (excerpt)}]
import threading
import customtkinter as ctk
import tkinter.messagebox as messagebox
from client import SecureClient
from crypto.hashing import DeepHash

ctk.set_appearance_mode("System")
ctk.set_default_color_theme("blue")

class SecureApp(ctk.CTk):
    def __init__(self):
        super().__init__()
        self.title("Secure Communication Suite")
        self.geometry("900x600")
        self.client = SecureClient()
        self.receive_thread = None
        self.running = True
        self.show_auth_frame()

    def authenticate(self, action):
        user = self.username_entry.get()
        pwd = self.password_entry.get()
        if not user or not pwd:
            messagebox.showwarning("Input Error",
                "Provide both username and password")
            return
        success, msg = self.client.connect_and_auth(
            action, user, pwd)
        if success:
            self.show_main_interface()
            self.start_receiving()
        else:
            messagebox.showerror("Auth Error", msg)

    def send_message(self):
        msg = self.msg_entry.get()
        if msg:
            success = self.client.send_message(msg)
            if success:
                self.msg_entry.delete(0, 'end')
                self.append_chat(f"You: {msg}")
            else:
                messagebox.showerror("Error",
                    "Failed to send. Lost connection.")
                self.disconnect()

    def receive_loop(self):
        while self.running and self.client.connected:
            msg = self.client.receive_message()
            if msg:
                self.after(0, self.append_chat, msg)
            else:
                if self.running:
                    self.after(0, self.disconnect)
                break

if __name__ == "__main__":
    app = SecureApp()
    app.mainloop()
\end{lstlisting}

% -----------------------------------------------------------
\subsection{How the Suite Secures Internet Services}
% -----------------------------------------------------------

The Secure Communication Suite applies its cryptographic modules to provide defence-in-depth for network communication:

\begin{enumerate}
    \item \textbf{Data in Transit:} All messages exchanged between client and server are encrypted using AES-EAX with per-session keys. This prevents eavesdropping (confidentiality) and detects any modification of packets (integrity).
    
    \item \textbf{Key Exchange Security:} The AES session key is never transmitted in plaintext. It is encrypted with the server's RSA-2048 public key, ensuring that only the legitimate server can recover it.
    
    \item \textbf{Data at Rest:} User credentials are stored as salted SHA-256 hashes, protecting them even if the database file is exfiltrated. RSA private keys are stored on disk in PEM format, with access controlled by file-system permissions.
    
    \item \textbf{Authentication:} Users must register and log in before accessing any service. The server rejects duplicate sessions, preventing impersonation or session hijacking.
    
    \item \textbf{Per-Client Isolation:} Each client has its own AES session key and cipher instance. A compromise of one client's key does not affect other clients' sessions.
\end{enumerate}

% -----------------------------------------------------------
\subsection{Testing}
% -----------------------------------------------------------

\subsubsection{Types of Tests Conducted}

\begin{table}[H]
\centering
\caption{Test plan summary}
\begin{tabular}{@{}p{3cm}p{4.5cm}p{5.5cm}@{}}
\toprule
\textbf{Test Type} & \textbf{Scope} & \textbf{Description} \\
\midrule
Unit Testing & Individual modules & Encrypt/decrypt round-trip, hash verification, register/login, key generation. \\
\addlinespace
Integration Testing & Module interaction & Full handshake: auth $\to$ key exchange $\to$ encrypted messaging. \\
\addlinespace
Functional Testing & End-to-end & Multi-client chat: messages sent by one client are correctly received and displayed by others. \\
\addlinespace
Security Testing & Cryptographic correctness & Tampered ciphertext raises exceptions; wrong passwords are rejected; duplicate logins are blocked. \\
\addlinespace
GUI Testing & User interface & Login/Register flow, message send/receive, hash computation, AES encrypt/decrypt tool. \\
\bottomrule
\end{tabular}
\end{table}

\subsubsection{Sample Test Scenarios}

\begin{enumerate}
    \item \textbf{AES Round-Trip:} Encrypt a plaintext message, then decrypt the result; verify the recovered plaintext matches the original.
    \item \textbf{Tamper Detection:} Flip a bit in the ciphertext and attempt decryption; verify that a \texttt{ValueError} is raised due to tag mismatch.
    \item \textbf{RSA Key Exchange:} Generate a key pair, encrypt a 16-byte AES key with the public key, decrypt with the private key, and verify equality.
    \item \textbf{Password Salting:} Register two users with the same password; verify that the stored hashes differ (because the salts differ).
    \item \textbf{Duplicate Login Prevention:} Log in as a user, then attempt a second login with the same username from a different client; verify the server returns \texttt{DUP}.
    \item \textbf{Multi-Client Broadcast:} Connect three clients, send a message from one, and verify that the other two receive it correctly decrypted.
\end{enumerate}

\subsubsection{Phase 4 Summary}
By the end of Phase~4, the Secure Communication Suite is a fully integrated application that:
\begin{itemize}
    \item Authenticates users via password-based credentials.
    \item Establishes secure sessions using RSA key exchange.
    \item Encrypts all communication with AES-EAX authenticated encryption.
    \item Provides SHA-256 hashing for independent data integrity verification.
    \item Offers a modern graphical interface for user-friendly interaction.
\end{itemize}

% ===========================================================
%  CONCLUSION
% ===========================================================
\newpage
\section{Conclusion}

The Secure Communication Suite successfully demonstrates the practical application of fundamental cryptographic techniques in a real-world Python application. Across three development phases (Phases~2--4), the project incrementally built, integrated, and tested six core security modules:

\begin{itemize}
    \item \textbf{AES-EAX} for authenticated symmetric encryption,
    \item \textbf{RSA-OAEP} for asymmetric key exchange,
    \item \textbf{SHA-256} for data integrity verification,
    \item \textbf{Key Management} for secure key lifecycle handling,
    \item \textbf{Password Authentication} with salted hashing,
    \item \textbf{Internet Services Security} via a multi-threaded client--server architecture.
\end{itemize}

The hybrid encryption approach (RSA for key transport, AES for bulk encryption) mirrors real-world protocols such as TLS, providing students with hands-on experience with industry-standard security practices.

% ===========================================================
%  REFERENCES
% ===========================================================
\section*{References}
\begin{enumerate}[label={[\arabic*]}]
    \item NIST, ``Advanced Encryption Standard (AES),'' FIPS PUB 197, 2001.
    \item R.~Rivest, A.~Shamir, and L.~Adleman, ``A Method for Obtaining Digital Signatures and Public-Key Cryptosystems,'' \textit{Communications of the ACM}, vol.~21, no.~2, pp.~120--126, 1978.
    \item NIST, ``Secure Hash Standard (SHS),'' FIPS PUB 180-4, 2015.
    \item M.~Bellare, P.~Rogaway, and D.~Wagner, ``The EAX Mode of Operation,'' \textit{FSE 2004}, LNCS 3017, pp.~389--407, 2004.
    \item PyCryptodome Documentation, \url{https://www.pycryptodome.org/}.
    \item CustomTkinter Documentation, \url{https://customtkinter.tomschimansky.com/}.
\end{enumerate}

\end{document}
